\documentclass[12pt,letterpaper]{article}
\usepackage{graphicx,textcomp}
\usepackage{natbib}
\usepackage{setspace}
\usepackage{fullpage}
\usepackage{color}
\usepackage[reqno]{amsmath}
\usepackage{amsthm}
\usepackage{fancyvrb}
\usepackage{amssymb,enumerate}
\usepackage[all]{xy}
\usepackage{endnotes}
\usepackage{lscape}
\newtheorem{com}{Comment}
\usepackage{float}
\usepackage{hyperref}
\newtheorem{lem} {Lemma}
\newtheorem{prop}{Proposition}
\newtheorem{thm}{Theorem}
\newtheorem{defn}{Definition}
\newtheorem{cor}{Corollary}
\newtheorem{obs}{Observation}
\usepackage[compact]{titlesec}
\usepackage{dcolumn}
\usepackage{tikz}
\usetikzlibrary{arrows}
\usepackage{multirow}
\usepackage{xcolor}
\newcolumntype{.}{D{.}{.}{-1}}
\newcolumntype{d}[1]{D{.}{.}{#1}}
\definecolor{light-gray}{gray}{0.65}
\usepackage{url}
\usepackage{listings}
\usepackage{color}

\definecolor{codegreen}{rgb}{0,0.6,0}
\definecolor{codegray}{rgb}{0.5,0.5,0.5}
\definecolor{codepurple}{rgb}{0.58,0,0.82}
\definecolor{backcolour}{rgb}{0.95,0.95,0.92}

\lstdefinestyle{mystyle}{
	backgroundcolor=\color{backcolour},   
	commentstyle=\color{codegreen},
	keywordstyle=\color{magenta},
	numberstyle=\tiny\color{codegray},
	stringstyle=\color{codepurple},
	basicstyle=\footnotesize,
	breakatwhitespace=false,         
	breaklines=true,                 
	captionpos=b,                    
	keepspaces=true,                 
	numbers=left,                    
	numbersep=5pt,                  
	showspaces=false,                
	showstringspaces=false,
	showtabs=false,                  
	tabsize=2
}
\lstset{style=mystyle}
\newcommand{\Sref}[1]{Section~\ref{#1}}
\newtheorem{hyp}{Hypothesis}

\title{Problem Set 2}
\date{Due: October 23, 2025}
\author{Applied Stats/Quant Methods 1}

\begin{document}
	\maketitle
	\section*{Instructions}
\begin{itemize}
	\item Please show your work! You may lose points by simply writing in the answer. If the problem requires you to execute commands in \texttt{R}, please include the code you used to get your answers. Please also include the \texttt{.R} file that contains your code. If you are not sure if work needs to be shown for a particular problem, please ask.
	\item Your homework should be submitted electronically on GitHub.
	\item This problem set is due before 23:59 on Thursday October 23, 2025. No late assignments will be accepted.

\end{itemize}

	
	\vspace{.5cm}
	\section*{Question 1: Political Science}
		\vspace{.25cm}
	The following table was created using the data from a study run in a major Latin American city.\footnote{Fried, Lagunes, and Venkataramani (2010). ``Corruption and Inequality at the Crossroad: A Multimethod Study of Bribery and Discrimination in Latin America. \textit{Latin American Research Review}. 45 (1): 76-97.} As part of the experimental treatment in the study, one employee of the research team was chosen to make illegal left turns across traffic to draw the attention of the police officers on shift. Two employee drivers were upper class, two were lower class drivers, and the identity of the driver was randomly assigned per encounter. The researchers were interested in whether officers were more or less likely to solicit a bribe from drivers depending on their class (officers use phrases like, ``We can solve this the easy way'' to draw a bribe). The table below shows the resulting data.

\newpage
\begin{table}[h!]
	\centering
	\begin{tabular}{l | c c c }
		& Not Stopped & Bribe requested & Stopped/given warning \\
		\\[-1.8ex] 
		\hline \\[-1.8ex]
		Upper class & 14 & 6 & 7 \\
		Lower class & 7 & 7 & 1 \\
		\hline
	\end{tabular}
\end{table}

\begin{enumerate}
	
	\item [(a)]
	Calculate the $\chi^2$ test statistic by hand/manually (even better if you can do "by hand" in \texttt{R}).\\
	\vspace{1 cm}
	
	I am creating two matrixes with the observed and expected values respectively in order to compare the observed and expected values and calculate $\chi^2$. I create a matrix with 2 rows and 3 columns, assign names to the rows and columns, and assign the given values. To calculate expected values for each cell, I multiply the sum of the row with the sum of the column and divide it by the total amount of observations. To calculate $\chi^2$ I will use the following formula: $\chi^2 = \sum{\frac{(f_{observed}-f_{expected})^2}{f_{expected}}}$
	\begin{verbatim}
		
		data <- matrix(NA, nrow = 2, ncol = 3)
		rownames(data) <- c("Upper class", "Lower Class")
		colnames(data) <- c("Not stopped", "Bribe requested", "Stopped/given warning")
		data[1,] <- c(14, 6, 7)
		data[2,] <- c(7, 7, 1)
		CSum <- colSums(data) 
		RSum <- rowSums(data)
		grand_total <- sum(data)
		
		expected <- outer(RSum, CSum)/grand_total
		
		chi_square <- sum((data - expected)^2/expected)
		
	\end{verbatim}
	
\textit{	Interpretation: The value of $\chi^2$ is 3.79. This is a low value, meaning that we will probably not be able to reject the upcoming H0.}
	
	
	\item [(b)]
	Now calculate the p-value from the test statistic you just created (in \texttt{R}).\footnote{Remember frequency should be $>$ 5 for all cells, but let's calculate the p-value here anyway.}  What do you conclude if $\alpha = 0.1$?\\
	
	Our null hypothesis s(H0) states that the variables are statistically independent. 
	
		\begin{verbatim}
		
		df_data <- (nrow(data)-1)*(ncol(data)-1)
		pchisq(chi_square, df = df_data, lower.tail = FALSE)
		
	\end{verbatim}
		
\textit{Interpretaion: 	The given p-Value = 0.15, which is larger than 0.1. As such, we do not reject H0 and the fact that variables are statistically independent. We can thus assume that the variables are statistically independent. }
	
	\item [(c)] Calculate the standardized residuals for each cell and put them in the table below.
	\vspace{1cm}
	
	To calculate adjusted standardized residuals, I will be applying:
	 $z =\frac{f_{observed}-f_{expected}}{se}$
	
	\begin{BVerbatim}
		
	se <- sqrt(expected*(1-RSum/grand_total)*(1-CSum/grand_total))
	z <- (data-expected)/se
	z
	
	\end{BVerbatim}

	\begin{table}[h]
		\centering
		\begin{tabular}{l | c c c }
			& Not Stopped & Bribe requested & Stopped/given warning \\
			\\[-1.8ex] 
			\hline \\[-1.8ex]
			Upper class  & 0.32  &   -1.51  & 1.64 \\
			\\
			Lower class &  -0.27 &  1.92 &   -1.52 \\
			
		\end{tabular}
	\end{table}
	
	
	\vspace{3cm}
	\item [(d)] How might the standardized residuals help you interpret the results?  
	
	 3 out of 6 cells have a higher value than expected (upper class/not stopped, lower class/Bribe requested, upper class/stopped), with the remaining 3 cells having a lower value than expected. However, the observed counts are not of great magnitude, they don't have a strong deviation. This aligns with our observation that the variables are statistically independent. 
	
	
\end{enumerate}
\newpage

\section*{Question 2: Economics}
Chattopadhyay and Duflo were interested in whether women promote different policies than men.\footnote{Chattopadhyay and Duflo. (2004). ``Women as Policy Makers: Evidence from a Randomized Policy Experiment in India. \textit{Econometrica}. 72 (5), 1409-1443.} Answering this question with observational data is pretty difficult due to potential confounding problems (e.g. the districts that choose female politicians are likely to systematically differ in other aspects too). Hence, they exploit a randomized policy experiment in India, where since the mid-1990s, $\frac{1}{3}$ of village council heads have been randomly reserved for women. A subset of the data from West Bengal can be found at the following link: \url{https://raw.githubusercontent.com/kosukeimai/qss/master/PREDICTION/women.csv}\\

\noindent Each observation in the data set represents a village and there are two villages associated with one GP (i.e. a level of government is called "GP"). Figure~\ref{fig:women_desc} below shows the names and descriptions of the variables in the dataset. The authors hypothesize that female politicians are more likely to support policies female voters want. Researchers found that more women complain about the quality of drinking water than men. You need to estimate the effect of the reservation policy on the number of new or repaired drinking water facilities in the villages.
\vspace{.5cm}
\begin{figure}[h!]
	\caption{\footnotesize{Names and description of variables from Chattopadhyay and Duflo (2004).}}
	\vspace{.5cm}
	\centering
	\label{fig:women_desc}
	\includegraphics[width=1.1\textwidth]{/Users/mairikachur/Documents/PhD/03. Classes/Quantitative Methods I/PS02/women_desc.png}
\end{figure}		

\newpage
\begin{enumerate}
	\item [(a)] State a null and alternative (two-tailed) hypothesis. 
	
	The authors hypothesize that female politicians are more likely to support policies female voters want, i.e. increasing the quality of drinking water. However, the variable "female" is not randomly assigned, as such we use the proxy "reservation". We want to assess if the reservation policy (i.e. assigning a treatment, a female village head) has an effect on the number of new or repaired drinking water facilities. If the treatment "reserved" had no effect, the mean of water facilities for "reserved = 0" and "reserved = 1" would be identical. 
	µ0 = mean number of new and repaired drinking water facilities when reserved = 0
	µ1 = mean number of new and repaired drinking water facilities when reserved = 1
	
	
	H0 = µ0 = µ1
	HA = µ0 $\neq$  µ1
	
	Our null hypothesis states that there is no difference in the mean number of new or repaired drinking water facilities between villages that received treatment ("reserved") and those that did not.
	
	\vspace{4 cm}
	\item [(b)] Run a bivariate regression to test this hypothesis in \texttt{R} (include your code!).
	
	Running a bivariate regression, regressing the outcome variable ("water") on the input variable ("reserved")
	
\begin{verbatim}
reg <- lm(water ~ reserved, data = data)
summary(reg)

Output: 

Coefficients:
                     Estimate Std. Error  t value Pr(>|t|)    
(Intercept)   14.738      2.286   6.446 4.22e-10 ***
reserved       9.252      3.948   2.344   0.0197 *  
---
Signif. codes:  0 ‘***’ 0.001 ‘**’ 0.01 ‘*’ 0.05 ‘.’ 0.1 ‘ ’ 1

Residual standard error: 33.45 on 320 degrees of freedom
Multiple R-squared:  0.01688,	Adjusted R-squared:  0.0138 
F-statistic: 5.493 on 1 and 320 DF,  p-value: 0.0197

\end{verbatim}
	
	
	
	\vspace{2cm}
	\item [(c)] Interpret the coefficient estimate for reservation policy. 
	
	The output above gives us an estimated coefficient for the intercept of 14.7, meaning that for x = 0, y = 14.7 on average. As such, $\beta_0$  = 14.7 is the expected number of new water facilities in villages that did not have reservations for women as council heads. 	
	
	$\beta_1$ = 9.2 is the additional expected amount of wells for villages where a female council head was reserved, as "reserved" holds the value X = 1. On average, we can expect villages with reserved female council heads to have 23.9 new or repaired drinking water facilities. 
	
	Additionally, the estimated coefficient for the variable "reserved" shows a p-value of 0.0197 and comes with a one-star (*) significance level, meaning that we can reject the null hypothesis that council head reservations don't have an effect on the amount of new or repaired drinking water facilities. We can interpret the result as the reservation policy having a positive effect on the amount of repaired and new drinking water facilities.
	 
\end{enumerate}

%\newpage
%	\section*{Question 3 (30 points): Biology}
%
%There is a physiological cost of reproduction for fruit flies, such that it reduces the lifespan of female fruit flies.  Is there a similar cost to male fruit flies?  This dataset contains observations from five groups of 25 male fruit flies. The experiment tests if increased reproduction reduces longevity for male fruit flies. The five groups are: males forced to live alone, males assigned to live with one or eight newly pregnant females (non-receptive females), and males assigned to live with one or eight virgin females (interested females). The name of the data set is \texttt{fruitfly.csv}.\footnote{Partridge and Farquhar (1981).``Sexual Activity and the Lifespan of Male Fruitflies''. \textit{Nature}. 294, 580-581.}
%	\vspace{1cm}
%
%\begin{tabular}{r|l}
%	\texttt{No} & serial number (1-25) within each group of 25\\
%	\texttt{type} & Type of experimental assignment \\
%	& \hspace{0.1in} $1=$ no females  \\
%	& \hspace{0.1in} $2=$ 1 newly pregnant female \\
%	& \hspace{0.1in} $3=$ 8 newly pregnant females\\
%	& \hspace{0.1in} $4=$ 1 virgin female\\
%	& \hspace{0.1in} $5=$ 8 virgin females\\
%	\texttt{lifespan} & lifespan (days)\\
%	\texttt{thorax} & length of thorax (mm)\\
%	\texttt{sleep} & percentage of each day spent sleeping\\
%\end{tabular}
%	\vspace{1cm}
%\begin{enumerate}
%	
%	\item
%	Import the data set and obtain summary statistiscs and examine the distribution of the overall lifespan of the fruitflies.  
%
%\newpage
%	\item
%	Plot \texttt{lifespan} vs \texttt{thorax}. Does it look like there is a linear relationship? Provide the plot. What is the correlation coefficient between these two variables?
%		\vspace{6cm}
%	\item
%	Regress \texttt{lifespan} on \texttt{thorax}.  Interpret the slope of the fitted model.
%			\vspace{6cm}
%	\item
%	Test for a significant linear relationship between  \texttt{lifespan} and \texttt{thorax}. Provide and interpret your results of your test.
%	
%\newpage
%	\item
%	
%	Provide the 90\% confidence interval for the slope of the fitted model.
%	
%			\vspace{.5cm}
%	\begin{itemize}
%		\item
%		Use the formula of confidence interval.		\vspace{.5cm}
%		\item
%		Use the function  \texttt{confint()}  in \texttt{R} .
%	\end{itemize}
%			\vspace{6cm}
%	\item Use the \texttt{predict()} function in \texttt{R} to (1) predict an individual fruitfly's lifespan when \texttt{thorax}=0.8 and (2) the average \texttt{lifespan} of fruitflies when \texttt{thorax}=0.8 by the fitted model. This requires that you compute prediction and confidence intervals. What are the expected values of lifespan? What are the prediction and confidence intervals around the expected values? 
%	
%			\vspace{6cm}
%	\item	For a sequence of \texttt{thorax} values, draw a plot with their fitted values for \texttt{lifespan}, as well as the prediction intervals and confidence intervals.
%


%\end{enumerate}
\end{document}
